\documentclass{resume} % Utilise le style personnalisé resume.cls

\usepackage[T1]{fontenc}
\usepackage[utf8]{inputenc}
\PassOptionsToPackage{hyphens}{url}
\usepackage[hidelinks]{hyperref}
\usepackage{needspace}
\usepackage[left=0.75in,top=0.6in,right=0.75in,bottom=0.6in]{geometry} % Marges du document
\newcommand{\tab}[1]{\hspace{.2667\textwidth}\rlap{#1}}
\newcommand{\itab}[1]{\hspace{0em}\rlap{#1}}

% Lignes datées avec retrait automatique (aucun tableau requis)
\newcommand{\rowD}[2]{\par\begingroup\leftskip=6.2em \noindent\hskip-6.2em\makebox[6.2em][l]{{\em #1}}#2\par\endgroup}
\newcommand{\rowA}[3]{\par\begingroup\leftskip=14em \noindent\hskip-14em\makebox[6.2em][l]{{\em #1}}\makebox[7.8em][l]{#2}#3\par\endgroup}

\name{ANTOINE DESROSIERS} % Nom
\address{antoine.desrosiers.6@ulaval.ca \\ \href{https://antoinedesrosiers.github.io/}{antoinedesrosiers.github.io}}
\address{581-997-2809}

\begin{document}

%----------------------------------------------------------------------------------------
%	ÉTUDES
%----------------------------------------------------------------------------------------

\begin{rSection}{Études}

{\bf Maîtrise en science politique avec mémoire} \hfill {\em Obtention prévue avril 2028}
\\ {\em Directeur de mémoire : Marc André Bodet}
\\ Université Laval, Québec

{\bf Baccalauréat en science politique} \hfill {\em 2026}
\\ Université Laval, Québec

{\bf Diplôme d'études collégiales} \hfill {\em 2022}
\\ {\em Concentration en commerce}
\\ Collège Champlain-Saint-Lawrence, Québec

{\bf Diplôme d'études secondaires} \hfill {\em 2020}
\\ {\em Sports-Arts-Études}
\\ Cardinal-Roy, Québec

\end{rSection}

%----------------------------------------------------------------------------------------
%	COMPÉTENCES PARTICULIÈRES
%----------------------------------------------------------------------------------------

\Needspace{6\baselineskip}
\begin{rSection}{Compétences particulières}

{\bf Cours universitaires pertinents~:}

«~Devis de recherche~» \hfill {\em automne 2026} \\
«~Méthodologie qualitative~» \hfill {\em hiver 2025} \\
«~Féminisme et philosophie~» \hfill {\em automne 2025}

{\bf Langues parlées et écrites~:} français, anglais (cégep anglophone)

{\bf Informatique~:} Final Cut Pro (montage vidéo), Motion (effets spéciaux vidéo), Designer (photo), Markdown, Covidence, Excel, R, GitHub, Dropbox, Zotero

\end{rSection}

%----------------------------------------------------------------------------------------
%	EXPÉRIENCE PERTINENTE EN RECHERCHE
%----------------------------------------------------------------------------------------

\Needspace{6\baselineskip}
\begin{rSection}{Expériences pertinentes en recherche}

{\bf Projet de mémoire} \hfill {\em 2026-}
\\ Université Laval, Québec
\begin{itemize}
	\vspace{-0.1cm}\item Ethnographie sur la culture musicale émergente de la ville de Québec
	\vspace{-0.3cm}\item Politiques d’équité, de diversité et d’inclusion (EDI)
	\vspace{-0.3cm}\item Réflexivité avec ma positition en tant que gérant et musicien
\end{itemize}

{\bf Auxiliaire de recherche dans l'équipe d'Eve Bourgeois} \hfill {\em 2026-}
\\ Université Laval, Québec
\begin{itemize}
	\vspace{-0.1cm}\item Revue systématique sur l’adaptation aux changements climatiques
	\vspace{-0.3cm}\item Alimentation d'une base de données
	\vspace{-0.3cm}\item Codage et lecture d'articles scientifiques
\end{itemize}

{\bf Auxiliaire de recherche dans l'équipe de Marc André Bodet} \hfill {\em 2026-}
\\ Université Laval, Québec
\begin{itemize}
	\vspace{-0.1cm}\item Vie parlementaire canadienne
	\vspace{-0.3cm}\item Alimentation d'une base de données
	\vspace{-0.3cm}\item Surveillance d'examens
\end{itemize}

{\bf Auxiliaire de recherche dans l'équipe de Marie Brossier} \hfill {\em 2025}
\\ Université Laval, Québec
\begin{itemize}
	\vspace{-0.1cm}\item Élites parlementaires et consolidation démocratique~: profils, renouvellement et professionnalisation au Sénégal (1959-2022)
	\vspace{-0.3cm}\item Alimentation d'une base de données
	\vspace{-0.3cm}\item Analyse de documents biographiques
	\vspace{-0.3cm}\item Collecte de fiches socio-biographiques de députés sénégalais
\end{itemize}

{\bf Projet sur la culture émergente de la ville de Québec} \hfill {\em 2025}
\begin{itemize}
	\vspace{-0.1cm}\item {\href{https://youtu.be/sqIpJmnHXpI}{Dans le cadre du cours {\em Méthode qualitative}, j'ai mené des entrevues avec des personnes importantes du milieu (artiste, propriétaire, organisateur) et leur ai remis mes conclusions par l'entremise d'une capsule vidéo professionnelle.}}
\end{itemize}

\end{rSection}

%----------------------------------------------------------------------------------------
%	AFFILIATIONS ACADÉMIQUES
%----------------------------------------------------------------------------------------

\Needspace{6\baselineskip}
\begin{rSection}{Affiliations académiques}

{\bf Centre pour l'étude de la citoyenneté démocratique (CÉCD)} \hfill {\em Depuis 2026}
\\ Québec

{\bf Chaire de recherche sur la démocratie et les institutions parlementaires} \hfill {\em Depuis 2026}
\\ Université Laval, Québec

\end{rSection}

%----------------------------------------------------------------------------------------
%	BOURSE ET DISTINCTION
%----------------------------------------------------------------------------------------

\Needspace{6\baselineskip}
\begin{rSection}{Bourses et distinctions}

{\bf Volet Pro de l'Ampli de Québec} \hfill {\em 2026-2027}
\\ Flore \& JADE, Québec
\begin{itemize}
	\vspace{-0.1cm}\item Bourse de 2 500 \$
    \vspace{-0.3cm}\item Accompagnement avec un mentor pendant 1 an pour la gestion de carrière
    \vspace{-0.3cm}\item Programme offert aux artistes émergents de la ville de Québec qui se démarquent par leur démarche artistique, la qualité de leur proposition sonore et visuelle ainsi que leur niveau de professionnalisation
\end{itemize}

{\bf Porteur de la masse} \hfill {\em 2026}
\\ Cérémonie de collation des grades, Faculté des sciences sociales, Université Laval, Québec
\begin{itemize}
	\vspace{-0.1cm}\item {\href{https://nouvelles.ulaval.ca/2026/06/26/antoine-desrosiers-le-politologue-musicien-6d584cde-e412-4fea-b561-124c3fdf04ae}{Un article a été publié sur mon parcours}}
\end{itemize}

{\bf Subvention Première Ovation - Musique} \hfill {\em 2026}
\\ Flore \& JADE, Québec
\begin{itemize}
	\vspace{-0.1cm}\item Subvention de 3 000 \$
    \vspace{-0.3cm}\item Enregistrement d'une première bande sonore pour un microalbum
\end{itemize}

{\bf Tableau d'honneur de la Faculté des sciences sociales} \hfill {\em 2025-2026}
\\ Faculté des sciences sociales, Université Laval, Québec
\begin{itemize}
	\vspace{-0.1cm}\item «~Cette distinction reconnaît les étudiantes et étudiants de chaque cohorte qui s'illustrent particulièrement par la qualité et la constance de leurs résultats et vous en faites partie grâce à votre excellente moyenne de cheminement de 3,92.~»
\end{itemize}

{\bf Mention d'honneur par la doyenne Aurélie Campana} \hfill {\em 2024-2025}
\\ Faculté des sciences sociales, Université Laval, Québec
\\ Critères~:
\begin{itemize}
	\vspace{-0.1cm}\item avoir maintenu une moyenne de programme égale ou supérieure à 3,50 sur 4,33
	\vspace{-0.3cm}\item se classer parmi les 5~\% des meilleurs dossiers du programme
\end{itemize}

{\bf Tableau d'honneur du Collège Champlain-Saint-Lawrence} \hfill {\em 2021-2022}
\\ Collège Champlain-Saint-Lawrence, Québec

\end{rSection}

%----------------------------------------------------------------------------------------
%	AUTRES EXPÉRIENCES
%----------------------------------------------------------------------------------------

\Needspace{6\baselineskip}
\begin{rSection}{Autres expériences}

{\bf Agent de recherche et de planification socio-économique} \hfill {\em Depuis 2024}
\\ Ministère de la Langue française, Québec

{\bf Moniteur en loisirs en CHSLD} \hfill {\em 2023}
\\ CHSLD Louis-Hébert, Québec

{\bf Service et entretien} \hfill {\em 2020-2023}
\\ Le Nordique Spa Stoneham, Stoneham et Tewkesbury

\end{rSection}

%----------------------------------------------------------------------------------------
%	ENGAGEMENTS SOCIAUX
%----------------------------------------------------------------------------------------

\Needspace{6\baselineskip}
\begin{rSection}{Engagements sociaux}

{\bf Musicien accompli «~Flore \& JADE~»} \hfill {\em Depuis 2022}
\begin{itemize}
	\vspace{-0.1cm}\item Groupe de musique pop-jazz
	\vspace{-0.3cm}\item Composition originale, écriture de paroles, coordination artistique et gérance
	\vspace{-0.3cm}\item Sortie d'un premier microalbum intitulé {\em Mauvais Mots, Mauvais Moments} enregistré dans le studio du {\em Pantoum} en 2025
	\vspace{-0.3cm}\item Plus d'une trantaine de réprésentations musicales devant public
    \vspace{-0.3cm}\item Participation a des concours majeurs, dont le {\em Festival International de la Chanson de Granby} et {\em Ma Première Place des Arts}
\end{itemize}

{\bf Création de {\em Journal Tévé}, booking et médias sociaux} \hfill {\em Depuis 2025}
\begin{itemize}
	\vspace{-0.1cm}\item Firme de production de contenu numérique (500~000+ impressions)
	\vspace{-0.3cm}\item Documente l'évolution de la vitalité artistique de la culture musicale émergente québécoise
	\vspace{-0.3cm}\item Collaboration avec des institutions importantes, dont le {\em Festival d'été de Québec, Festival FONO, Le Pantoum} et {\em le Phoque Off}
	\vspace{-0.3cm}\item Collaboration avec des artistes à succès
	\vspace{-0.1cm}\item Maire de Laval, Kinji00, Marie Céleste, Velours Velours
\end{itemize}

{\bf Club Entrepreneuriat Saint-Lawrence (cégep)} \hfill {\em 2021-2022}
\begin{itemize}
	\vspace{-0.1cm}\item Comité marketing et réseaux sociaux
	\vspace{-0.3cm}\item Organisation d'événements socioculturels
\end{itemize}

\end{rSection}

%----------------------------------------------------------------------------------------
%	REPRÉSENTATIONS MUSICALES
%----------------------------------------------------------------------------------------

\Needspace{6\baselineskip}
\begin{rSection}{Représentations musicales}

{\bf 35 représentations musicales}

{\small\parskip=2pt
    \smallskip
	\rowA{2026/11/18}{Flore \& JADE}{Quart de finale à la 33e édition du concours Ma Première Place des Arts}
	\rowA{2026/10/30}{Flore \& JADE}{Première partie de Pierre Flynn -- Chez Madame Belley}
	\rowA{2026/09/09}{Flore \& JADE}{Show de la rentrée à Université Laval}
	\rowA{2026/08/29}{Flore \& JADE}{Saint-Roch-XP (2)}
	\rowA{2026/08/29}{Flore \& JADE}{Saint-Roch-XP (1)}
	\rowA{2026/08/19}{Flore \& JADE}{Demi-finaliste au Grand Concours Hydro-Québec -- Festival International de la Chanson de Granby 2026 (FICG)}
	\rowA{2026/08/15}{Flore \& JADE}{Auberge des Balcons à Baie-Saint-Paul}
	\rowA{2026/08/13}{Flore \& JADE}{Soirée : Chanson imposée (FICG)}
	\rowA{2026/08/11}{Flore \& JADE}{Les jardins du bonheur (FICG)}
	\rowA{2026/07/31}{Flore \& JADE}{Café Zénob à Trois-Rivières}
	\rowA{2026/07/17}{Flore \& JADE}{Animation Baie-Saint-Paul (3)}
	\rowA{2026/06/13}{Flore \& JADE}{Fête de la musique du Québec 19e édition}
	\rowA{2026/05/28}{Flore \& JADE}{Lancement du EP~: «~Mauvais Mots, Mauvais Moments~»}
	\rowA{2025/11/20}{Flore \& JADE}{Série des membres de l'Ampli de Québec}
	\rowA{2025/07/19}{Flore \& JADE}{Ambiance au Quai 1635 (2)}
	\rowA{2025/07/03}{Flore \& JADE}{Ambiance au Quai 1635 (1)}
	\rowA{2025/06/27}{Flore \& JADE}{Hotel le Capitaine à Isle-Aux-Coudres}
	\rowA{2025/06/21}{Flore \& JADE}{Fête de la musique du Québec 18e édition}
	\rowA{2025/06/19}{Flore \& JADE}{Université Laval en Lumière}
	\rowA{2025/04/12}{Flore \& JADE}{Première partie du groupe Bonace}
	\rowA{2025/02/28}{Flore \& JADE}{Spectacle à CHYZ}
	\rowA{2025/01/17}{Flore \& JADE}{Production culturelle}
	\rowA{2024/12/20}{Flore \& JADE}{Plateau double avec Maman}
	\rowA{2024/12/05}{Flore \& JADE}{Attaboy Pizzeria}
	\rowA{2024/11/24}{Flore \& JADE}{Marché de Noël allemand}
	\rowA{2024/10/26}{Flore \& JADE}{Lancement du projet au Fou-Bar}
	\rowA{2024/06/16}{Flore \& JADE}{Fête de la musique du Québec 17e édition}
	\rowA{2024/05/11}{Flore \& JADE}{Spectacle d'anniversaire à Providence, RI}
	\rowA{2023/08/19}{Flore \& JADE}{Animation Baie-Saint-Paul (2)}
	\rowA{2023/05/12}{Flore \& JADE}{«~Créations Printanières~»}
	\rowA{2023/04/08}{Flore \& JADE}{Pré-lancement de projet à La Librairie St-Jean-Baptiste}
	\rowA{2022/12/11}{Perles}{Marché de Noël (2)}
	\rowA{2022/11/10}{Perles}{Marché de Noël (1)}
	\rowA{2022/10/16}{Perles}{Animation Baie-Saint-Paul (1)}
	\rowA{2022/09/10}{Perles}{Lancement de projet}
}

\end{rSection}

%----------------------------------------------------------------------------------------
%	ENREGISTREMENTS VIDÉO
%----------------------------------------------------------------------------------------

\Needspace{6\baselineskip}
\begin{rSection}{Enregistrements vidéo}

{\bf À la réalisation et au montage}

{\small\parskip=2pt
    \smallskip
	\rowD{2026/07/28}{\href{https://youtu.be/I8BqAwUyDP8}{Journal Tévé -- Le lancement de Flore \& JADE s'est bien déroulé\dots} \\ {\em Réalisateur, monteur}}
	\rowD{2026/02/09}{\href{https://youtu.be/S5CMl92JqWg}{Journal Tévé~: Un vrai reportage sur le Maire de Laval (réalisé par JTV)} \\ {\em Réalisation, entrevue, montage}}
	\rowD{2025/07/31}{\href{https://youtu.be/-JRZWwhhqUo}{Journal Tévé -- Un genre de podcast avec MARIE CÉLESTE} \\ {\em Réalisateur, entrevue, montage}}
	\rowD{2025/04/23}{\href{https://youtu.be/sqIpJmnHXpI}{LA CAPITALE NATIONALE EN DÉCLIN~?} \\ {\em Chercheur, entrevue, montage, réalisation}}
	\rowD{2024/02/04}{\href{https://youtu.be/gQD-LTMbfkU}{«~Créations Printanières~» Flore \& JADE -- Stoneham 05/12/2023} \\ {\em Création et batteur}}
	\rowD{2023/03/30}{\href{https://youtu.be/zPtooWwGqF8}{deroza -- Accident (Official video) (Directed by Tristan TA)} \\ {\em Création et chanteur}}
}

\end{rSection}

%----------------------------------------------------------------------------------------
%	PROPRIÉTÉ INTELLECTUELLE
%----------------------------------------------------------------------------------------

\Needspace{6\baselineskip}
\begin{rSection}{Propriété intellectuelle}

{\bf Droits d'auteur enregistrés} \\ Composition, arrangement et production -- SOCAN. Année d'émission~: 2026.

{\small\parskip=2pt
	\smallskip
	\rowD{2026/11/10}{Et je m'élance -- Flore \& JADE \\ {\em Financement~: Première Ovation - Musique}}
	\rowD{2026/11/10}{Histoire de famille \\ {\em Financement~: Première Ovation -- Musique}}
	\rowD{2026/11/10}{Café Klem -- Flore \& JADE \\ {\em Financement~: Première Ovation -- Musique}}
	\rowD{2026/11/10}{Je sais, je sais -- Flore \& JADE \\ {\em Financement~: Première Ovation -- Musique}}
}

\end{rSection}

%----------------------------------------------------------------------------------------
%	ÉMISSIONS DE RADIO ET ENTREVUES PUBLIÉES
%----------------------------------------------------------------------------------------

\Needspace{8\baselineskip}
\begin{rSection}{Émissions de radio et entrevues publiées}

{\bf Émissions de radio et de télévision}

{\small\parskip=2pt
	\smallskip
    \rowD{2026/11}{{\em Chéri.e, j'arrive~!} \\ {\em CHYZ 94,3 FM}}
	\rowD{2026/08}{{\em Le Téléjournal de Radio-Canada}~ \\ {\em ICI Radio-Canada}}
	\rowD{2026/08}{{\em Première heure avec Alexandre Duval}~ \\ {\em ICI Radio-Canada Radio}}
	\rowD{2026/07}{{\em Les matins éphémères}~ \\ {\em CKRL 89,1}}
	\rowD{2026/05}{{\em Première heure avec Alexandre Duval}~ \\ {\em ICI Radio-Canada Radio}}
	\medskip
}

{\bf Entrevues publiées}

{\small\parskip=2pt
	\smallskip
	\rowD{2026/09/14}{\href{https://impactcampus.ca/arts-et-culture/le-tout-debut-de-lautomne-retour-sur-le-show-de-la-rentree/}{Flore \& JADE au Show de la Rentrée de l'Université Laval 2026} \\ {\em Impact Campus}}
	\rowD{2026/05/23}{Flore \& JADE lance un premier microalbum aux sonorités pop-jazz éclectiques \\ {\em Le Charlevoisien}}
	\rowD{2024/11/03}{Flore \& JADE, pop jazz née dans Charlevoix \\ {\em Le Charlevoisien}}
}

\end{rSection}

\end{document}